% s05-grassmannian.tex —— §3 的 B_dR-仿射 Grassmann 簇与 Hecke 叠（原文 p.23–p.28）
% 底本：L. Fargues, Geometrization of the local Langlands correspondence: an overview, 2016-02-02
% 体例：每页起点 \srcpage{N}；定理/命题编号与原文逐条一致（共享 [section] 计数器）；
%       译注 \yizhu{注文}{来源}，无外部依据用 \yizhuw{}。
% 注：本文件不写 documentclass / \begin{document}。

% 原文用 ⟦⟧ 表示幂级数环；preamble 未载 stmaryrd，故就地兜底（s02 已有同款定义，重复无害）。

\section{\texorpdfstring{$B_{\mathrm{dR}}$}{BdR}-仿射 Grassmann 簇与 Hecke 叠}% 原文 3. The BdR-affine Grassmanian and the Hecke stack
\srcpage{23}
\subsection{\texorpdfstring{$B_{\mathrm{dR}}$}{BdR}-仿射 Grassmann 簇：定义}% 原文 3.1. The BdR-affine Grassmannian: definition

\leavevmode 当 $E=\FF_q((\pi))$ 时，我们采用约定 $\Spa(E)^{\diamond}=\Spa(E)$\yizhu{当 $E=\FF_q((\pi))$ 时 $\Spa(E)$ 本身已是完满空间，故其菱形就是自身；对一般的 $E$，$(-)^{\diamond}$ 表示取完满化。}{A：原文 §1.3.1 p.7；I：本工程 glossary.tsv（diamond）}。$B_{\mathrm{dR}}$-仿射 Grassmann 簇按定义是一个 pro-étale 层\yizhuw{原书此处以竖箭头排版，示意 $\Gr$ 是 $\Spa(E)^{\diamond}$ 上的层，并非一般意义的态射图。译文照原样排出。}：

\[
\begin{tikzcd}
  \Gr \arrow[d] \\
  \Spa(E)^{\diamond}
\end{tikzcd}
\]

也就是说，它是 $\mathrm{Perf}_E$ 上的一个 pro-étale 层，具体定义如下。

\begin{definition}\label{def:3.1}
设 $(R,R^+)$ 是 $E$ 上的 affinoid 完满代数。则
\[
  \Gr(R,R^+)=\{(\mathcal{T},\xi)\}/\!\sim,
\]
其中 $\mathcal{T}$ 是 $\Spec(B_{\mathrm{dR}}^+(R))$\yizhu{$B_{\mathrm{dR}}^+(R)$ 是 Fontaine 的 de Rham 周期环的正部分。当 $E=\FF_q((\pi))$ 时它同构于 $R\llbracket T\rrbracket$（见 §3.3.1）；当 $E|\QQ_p$ 时由 §1.6.2 的构造给出。}{A：原文 §1.6.2 p.12；Fontaine, \emph{Représentations p-adiques}（1979）} 上的一个 $G$-扭子，$\xi$ 是 $\mathcal{T}\otimes_{B_{\mathrm{dR}}^+(R)}B_{\mathrm{dR}}(R)$ 的一个平凡化。
\end{definition}

\srcpage{24}
\begin{remark}\label{rem:3.2}
按定义，$\Gr(R,R^+)$ 只依赖 $R$，不依赖 $R^+$ 的选取。$\Bun_G(R,R^+)$、$\Spa(E)^{\diamond}$，以及文中出现的所有层与叠，都有同样的现象。从几何上看，这说明我们所考虑的一切几何对象都是\emph{部分 proper}（partially proper）的：用 Tate 刚性几何的话说，就是超收敛的；用 Berkovich 的话说，就是无边界的。
\end{remark}

\begin{remark}\label{rem:3.3}
读者也许会奇怪：$\Gr$ 的定义里出现的是“Spec”，而不是“Spa”。只要想到，对任意 affinoid 完满代数 $(R,R^+)$（Kedlaya–Liu）有等价
\[
  \text{$\Spa(R,R^+)$ 上的向量丛}\;\simeq\;\text{$\Spec(R)$ 上的向量丛},
\]
就不会再觉得奇怪。于是 $\Spa(R,R^+)$ 上 Tannakian 意义下的 $G$-丛，与 $\Spec(R)$ 上按 étale 拓扑局部平凡的 $G$-扭子是一回事。
\end{remark}

\leavevmode 可以验证：$\Spa(R,R^+)$ 上的任何 $G$-丛，经过 $\Spa(R,R^+)$ 的一个 étale 覆盖后都变为平凡。由此推出，$\Gr$ 就是预层
\[
  (R,R^+)\longmapsto G(B_{\mathrm{dR}}(R))/G(B_{\mathrm{dR}}^+(R))
\]
所对应的 étale 层。

\subsection{Beauville–Laszlo 单值化态射}% 原文 3.2. The Beauville-Laszlo uniformization morphism

设 $(R,R^+)$ 是 $\FF_q$ 上的 affinoid 完满代数，并考虑 $R$ 在 $E$ 上的一个 untilt $R^\sharp$。这就定义了相对概形曲线 $X_R^{sch}$ 上的一个 Cartier 除子。沿该除子的形式完备化是 $\Spf(B_{\mathrm{dR}}^+(R^\sharp))$。应用 Beauville–Laszlo 黏合（[1]），我们得到以下结论。

\begin{proposition}\label{prop:3.4}
作为 $\Spa(E)^{\diamond}$ 上的层，
\[
  \Gr(R,R^+)=\{(R^\sharp,\mathcal{E},\xi)\}/\!\sim,
\]
其中
\begin{itemize}
  \item $R^\sharp$ 是 $R$ 的一个 untilt；
  \item $\mathcal{E}$ 是 $X_R^{sch}$ 上的 $G$-丛；
  \item 若 $D$ 是 $X_R^{sch}$ 上由 $R^\sharp$ 定义的 Cartier 除子，则 $\xi$ 是 $\mathcal{E}|_{X_R^{sch}\setminus D}$ 的一个平凡化。
\end{itemize}
\end{proposition}

作为推论，有 Beauville–Laszlo 态射
\[
  \mathcal{BL}:\Gr\longrightarrow\Bun_G\times_{\Spa(\FF_q)}\Spa(E)^{\diamond}.
\]

回忆下面这条结果，它是 Drinfeld–Simpson 一个结论的类比（[5]）。

\begin{proposition}[{[6]}]\label{prop:3.5}
设 $F|\FF_q$ 是完满的代数闭域，$\infty\in|X_F^{sch}|$ 是闭点。则 $X_F^{sch}$ 上的任何 $G$-丛在 $X_F^{sch}\setminus\{\infty\}$ 上都变为平凡。
\end{proposition}

\begin{remark}\label{rem:3.6}
Drinfeld 与 Simpson 附加了 $G$ 半单这一假设。在我们的情形不需要，因为对这样的 $F$ 有 $\Pic^0(X_F)=0$。
\end{remark}

下面是前一命题的几何表述。

\begin{proposition}\label{prop:3.7}
Beauville–Laszlo 态射 $\mathcal{BL}$ 在几何点层面是满射。
\end{proposition}

我们自然猜想如下。

\begin{conjecture}\label{conj:3.8}
态射 $\mathcal{BL}$“在光滑拓扑下局部满”。更精确地说，对 $S\in\mathrm{Perf}_{\FF_q}$ 及 $(\Bun_G\times\Spa(E)^{\diamond})(S)$ 的一个元素，存在一个“光滑”满覆盖 $\widetilde{S}\to S$，使得该元素限制到 $\widetilde{S}$ 后来自 $\Gr(\widetilde{S})$ 的一个元素。
\end{conjecture}

仿照 Drinfeld–Simpson，这可由下面这条推出。

\begin{conjecture}\label{conj:3.9}
对 $S\in\mathrm{Perf}_{\FF_q}$ 与 $X_S$ 上的 $G$-丛 $\mathcal{E}$，$\mathcal{E}$ 到 $B$ 的约化的模空间是 $S$ 上的一个“光滑”菱形；当 $E=\FF_q((\pi))$ 时，则是完满空间。
\end{conjecture}

这应该可由 §2.1 中那样的 Quot-菱形技术推出。

\subsection{\texorpdfstring{$B_{\mathrm{dR}}$}{BdR}-仿射 Grassmann 簇上的几何结构}% 原文 3.3. Geometric structure on the BdR-affine Grassmanian

\srcpage{25}
\subsubsection*{3.3.1. 等特征情形}% 原文 3.3.1. The equal characteristic case.

\leavevmode 我们从 $E=\FF_q((\pi))$ 的情形开始。此时，读者现在应该已经明白，菱形不应出现，我们应当只需与完满空间打交道。

回忆（§1.6.2）：此时对完满 $E$-代数 $R$ 有 $B_{\mathrm{dR}}^+(R)=R\llbracket T\rrbracket$，其中结构态射 $E\to R\llbracket T\rrbracket$ 由
\[
  \pi\longmapsto T+a
\]
给出；这里 $\pi\mapsto a\in R^{\circ\circ}\cap R^\times$ 经由 $E\to R$ 得到，并且
\[
  T+a=a\Bigl(1+\frac{T}{a}\Bigr)\in R\llbracket T\rrbracket^\times .
\]

先设 $G$ 非分歧，并记 $\underline{G}$ 为它在 $\oo_E$ 上的约化模型。记 $\mathrm{Nilp}_{\oo_E}$ 为 $\pi$ 在其上幂零的 $\oo_E$-代数所成的范畴。我们定义函子
\[
  \mathcal{F}:\mathrm{Nilp}_{\oo_E}\longrightarrow\mathrm{Sets}.
\]
设 $R\in\mathrm{Nilp}_{\oo_E}$，$\pi\mapsto a\in R$。给 $R\llbracket T\rrbracket$ 装备由 $\pi\mapsto T+a$ 给出的 $\FF_q\llbracket\pi\rrbracket$-代数结构。由于 $a$ 幂零，$R((T))$ 是一个 $E$-代数。我们于是令
\[
  \mathcal{F}(R)=G\bigl(R((T))\bigr)/G(R\llbracket T\rrbracket).
\]
现记 $\widetilde{\mathcal{F}}$ 为 $\mathcal{F}$ 所对应的 étale 层。于是可以验证以下结论。

\begin{proposition}\label{prop:3.10}
\begin{enumerate}[label=(\arabic*)]
  \item étale 层 $\widetilde{\mathcal{F}}$ 可由一个 ind $\pi$-进 $\oo_E$-形式概形表示。该形式概形局部有限型，其特殊纤维
  \[
    \Gr^{\mathrm{PR}}:=\widetilde{\mathcal{F}}\otimes\oo_E/\pi
  \]
  是联系于某个超特殊子群的 Pappas–Rapoport 扭仿射 Grassmann 簇（[21]）。
  \item 有同构
  \[
    \widetilde{\mathcal{F}}\simeq\Gr^{\mathrm{PR}}\widehat{\otimes}_{\FF_q}\oo_E .
  \]
  \item 记 $\widetilde{\mathcal{F}}_\eta$ 为 $\widetilde{\mathcal{F}}$ 所伴随的 ind-刚性解析空间，即它的一般纤维。则
  \[
    \Gr=\varprojlim_{\mathrm{Frob}}\widetilde{\mathcal{F}}_\eta
  \]
  是一个 ind 完满空间，它是 $\widetilde{\mathcal{F}}_\eta$ 的完满化。
\end{enumerate}
\end{proposition}

第 (1) 点由第 (2) 点推出。第 (2) 点是下述事实的推论：对 $R\in\mathrm{Nilp}_{\oo_E}$ 且 $\pi\mapsto a\in R$，对应 $T\mapsto T-a$ 定义了 $R[T]$ 的一个自同构；因 $a$ 幂零，它连续地延拓为 $R\llbracket T\rrbracket$ 的自同构。

若 $Z$ 是 $\FF_q$ 上的 proper 概形，则 $(Z\widehat{\otimes}_{\FF_q}\oo_E)_\eta=(Z\otimes_{\FF_q}E)^{\mathrm{rig}}$。由于 Pappas–Rapoport 扭仿射 Grassmann 簇是 proper 概形的归纳极限，我们推出以下结论。

\begin{corollary}\label{cor:3.11}
设 $G$ 非分歧。$B_{\mathrm{dR}}$-仿射 Grassmann 簇 $\Gr$ 是一个 ind-完满空间。取联系于某个超特殊子群的 Pappas–Rapoport 扭仿射 Grassmann 簇，由 $\FF_q$ 标量扩张到 $E$，再作解析化并取完满化，所得就是它：
\[
  \Gr^{B_{\mathrm{dR}}}=\varprojlim_{\mathrm{Frob}}\bigl(\Gr^{\mathrm{PR}}\otimes_{\FF_q}E\bigr)^{\mathrm{rig}}.
\]
\end{corollary}

用从非分歧情形出发的下降理论，可以容易地处理一般情形。下面这个命题对 $G$ 没有任何假设。

\begin{proposition}\label{prop:3.12}
设 $R$ 是 affinoid $E$-代数，给 $R\llbracket T\rrbracket$ 装备由 $E\ni\pi\mapsto T+\pi$ 给出的 $E$-代数结构。
\begin{enumerate}[label=(\arabic*)]
  \item 预层 $R\mapsto G(R((T)))/G(R\llbracket T\rrbracket)$ 所对应的 étale 层可由一个 ind-$E$-刚性解析空间表示。
  \item 这个 ind-刚性解析空间的完满化是一个 ind-完满空间，它表示 $B_{\mathrm{dR}}$-仿射 Grassmann 簇。
\end{enumerate}
\end{proposition}

\srcpage{26}
\begin{remark}\label{rem:3.13}
须注意
\[
  \pi_0(\Gr_{B_{\mathrm{dR}}})=\pi_1(G)/\Gamma .
\]
当 $G$ 非分歧时，这与 [21] 一致；当 $G$ 分歧时，则与 [21] 不同。
\end{remark}

\subsubsection*{3.3.2. 不等特征情形}% 原文 3.3.2 印作 The equal characteristic case.（笔误，见译注）

\leavevmode 现设 $E|\QQ_p$。若 $C|E$ 完备且代数闭，则有双射\yizhuw{原文小节标题印作“The equal characteristic case”，但本小节通篇假设 $E|\QQ_p$，显系“unequal characteristic”之误。译文按内容改题，未改原文。}
\[
  X_*(T)^+\xrightarrow{\ \sim\ }G(B_{\mathrm{dR}}^+(C))\backslash G(B_{\mathrm{dR}}(C))/G(B_{\mathrm{dR}}^+(C)).
\]
这定义了一个相对位置映射
\[
  |\Gr|\longrightarrow X_*(T)^+/\Gamma,
\]
它半连续。于是可以写
\[
  \Gr=\varinjlim_{\mu\in X_*(T)^+/\Gamma}\Gr^{\le\mu}.
\]
于是有 Scholze 的下述定理。

\begin{theorem}[{[27]}]\label{thm:3.14}
对每个 $\mu\in X_*(T)^+/\Gamma$，pro-étale 层 $\Gr^{\le\mu}$ 是 $\Spa(E)^{\diamond}$ 上的菱形。
\end{theorem}

Scholze 给出的证明经由所谓的忠实拓扑，是非构造性的。尽管如此，人们可以期望：像 §2.1 那样使用 Quot-菱形技术，能为这个定理找到另一个更具构造性的证明（经由 $B_{\mathrm{dR}}$-仿射 Grassmann 簇与 Hecke 叠之间的联系）。

下面是受等特征情形启发的另一个期望\yizhuw{原文用“Hope”而不用“Conjecture”，语气弱于猜想：作者只表示期望这样的对象存在，并未提出明确猜想。preamble 已把该环境定名为“期望”。}。

\begin{hope}\label{hope:3.15}
若 $G$ 非分歧，则存在一个对象，它居于 $\Spa(\oo_E)^{\diamond}$ 之上：其在 $\FF_q$ 上的特殊纤维是 Zhu 所定义的 Witt 向量仿射 Grassmann 簇（[32]、[29]）\yizhu{Witt 向量仿射 Grassmann 簇是 $p$-进情形下与 $B_{\mathrm{dR}}$-仿射 Grassmann 簇平行的对象：以 Witt 向量环取代 $B_{\mathrm{dR}}^+$，住在 $\Spa(\oo_E)$ 的特殊纤维一侧。}{A：原文 p.26（引 [32]、[29]）}，其在 $E^{\diamond}$ 上的一般纤维是 $B_{\mathrm{dR}}$-仿射 Grassmann 簇。
\end{hope}

\begin{remark}\label{rem:3.16}
须小心：$\Spa(\oo_E)^{\diamond}$ 先验地没有几何结构；它只是一个 pro-étale 层，而不是菱形。
\end{remark}

\subsection{Hecke 叠}% 原文 3.4. The Hecke stack

固定 $\mu\in X_*(T)^+/\Gamma$，作如下定义。

\begin{definition}\label{def:3.17}
设 $(R,R^+)$ 是完满 affinoid $\FF_q$-代数。$\mathit{Hecke}^{\le\mu}(R,R^+)$ 是四元组 $(\mathcal{E}_1,\mathcal{E}_2,R^\sharp,f)$ 所成的群胚，其中：
\begin{itemize}
  \item $R^\sharp$ 是 $R$ 在 $E$ 上的一个 untilt；
  \item $\mathcal{E}_1$ 与 $\mathcal{E}_2$ 是 $X_R^{sch}$ 上的 $G$-丛；
  \item 若 $D$ 是 $X_R^{sch}$ 上由 untilt $R^\sharp$ 定义的 Cartier 除子，则
  \[
    f:\mathcal{E}_1|_{X_R^{sch}\setminus D}\xrightarrow{\ \sim\ }\mathcal{E}_2|_{X_R^{sch}\setminus D};
  \]
  \item 在 $\Spa(R,R^+)$ 上逐纤维地，修正 $f$ 以 $\mu$ 为界。
\end{itemize}
\end{definition}

我们把最后一个条件说得更精确。设 $\Spa(F)\to\Spa(R,R^+)$ 是几何点，untilt $R^\sharp$ 定义了 $F$ 的一个 untilt $C$，它给出闭点 $\infty\in X_F^{sch}$，并满足 $\widehat{\oo}_{X_F^{sch},\infty}=B_{\mathrm{dR}}^+(C)$。对 $i=1,2$，固定 $\mathcal{E}_i\times_{X_F^{sch}}\Spec(B_{\mathrm{dR}}^+(C))$ 的一个平凡化。于是 $f$ 拉回到 $\Spec(B_{\mathrm{dR}}(C))$ 上，由 $G(B_{\mathrm{dR}}(C))$ 的一个元素给出。模去上述平凡化的选取，这定义了
\[
  X_*(T)^+=G(B_{\mathrm{dR}}^+(C))\backslash G(B_{\mathrm{dR}}(C))/G(B_{\mathrm{dR}}^+(C))
\]
的一个元素。我们要求这个元素以 $\mu$ 的 Galois 轨道中的某个元素为界。

\srcpage{27}
\begin{remark}\label{rem:3.18}
我们本可以对任意 $S\in\mathrm{Perf}_{\FF_q}$ 直接定义 $\mathit{Hecke}^{\le\mu}(S)$，但那个定义更微妙。事实上必须使用 adic 曲线 $X_S^{ad}$；此时两个 $G$-丛之间修正的定义更复杂，因为必须要求 $X_S^{ad}\setminus D$ 上的同构“沿 $D$ 亚纯”。这里我们用到了概形曲线的一大好处：修正的定义更简单（而且 Beauville–Laszlo 黏合可直接适用于它；对 adic 曲线虽也可适用，但要稍微复杂一些）。
\end{remark}

\begin{definition}\label{def:3.19}
我们定义如下交换图

\[
\begin{tikzcd}
  & \mathit{Hecke}^{\le\mu} \arrow[dl, "\overleftarrow{h}"'] \arrow[dr, "\overrightarrow{h}"] & \\
  \Bun_G & & \Bun_G\times\Spa(E)^{\diamond}
\end{tikzcd}
\]

其中两个态射由公式
\[
  \overleftarrow{h}(\mathcal{E}_1,\mathcal{E}_2,R^\sharp,f)=\mathcal{E}_2
  \quad\text{与}\quad
  \overrightarrow{h}(\mathcal{E}_1,\mathcal{E}_2,R^\sharp,f)=(\mathcal{E}_1,R^\sharp)
\]
给出。
\end{definition}

定义 $L^+G$ 为 $E^{\diamond}$ 上的一个群，公式为
\[
  L^+G(R,R^+)=G(B_{\mathrm{dR}}^+(R)),
\]
其中 $(R,R^+)$ 是完满 affinoid $E$-代数。存在一个 étale $L^+G$-扭子
\[
  \mathcal{T}\longrightarrow\Bun_G\times\Spa(E)^{\diamond}.
\]
事实上，给定完满 $\FF_q$-代数 $R$ 及 $R$ 在 $E$ 上的 untilt $R^\sharp$，对 $X_R^{sch}$ 上的 $G$-丛 $\mathcal{E}$，它拉回到 $\Spec(B_{\mathrm{dR}}^+(R^\sharp))$ 上，就定义了该概形上的一个 $G$-扭子。

群 $L^+G$ 作用在 $B_{\mathrm{dR}}$-仿射 Grassmann 簇上。下面这个命题再次用到 Beauville–Laszlo 黏合。

\begin{proposition}\label{prop:3.20}
态射
\[
  \overrightarrow{h}:\mathit{Hecke}^{\le\mu}\longrightarrow\Bun_G\times\Spa(E)^{\diamond}
\]
等同于
\[
  \mathcal{T}\times_{L^+G}\Gr^{\le\mu}.
\]
换言之，$\overrightarrow{h}$ 是以 $\Gr^{\le\mu}$ 为纤维的 étale 局部平凡纤维化，其黏合态射由 $L^+G$ 中的元素给出。
\end{proposition}

\subsection{假想的几何 Satake 对应}% 原文 3.5. A hypothetical geometric Satake correspondence

固定 $\ell\neq p$。记 ${}^LG=\widehat{G}\rtimes\Gamma$ 为 $G$ 的 $\overline{\QQ}_\ell$-Langlands 对偶\yizhu{这里 $\Gamma=\Gal(\overline{E}|E)$ 通过外作用与 $\widehat{G}$ 作半直积；“$\overline{\QQ}_\ell$-对偶”指取 $\ell$-进表示意义下的对偶群，与 §4 的用法一致。}{A：原文 §3.5 p.27、§4 p.28}。

\begin{conjecture}\label{conj:3.21}
\begin{enumerate}[label=(\arabic*)]
  \item 存在 $\Gr$ 上一个 $L^+G$-等变的 $\overline{\QQ}_\ell$-“反常层”范畴，它等价于 $\Rep_{\overline{\QQ}_\ell}({}^LG)$，即 ${}^LG$ 在有限维 $\overline{\QQ}_\ell$-向量空间中的表示所成范畴；这些表示限制到 $\widehat{G}$ 上是代数的，且在 ${}^LG$ 的 $\Gamma$-因子上是离散的。
  \item 经由这个等价，对 $\mu\in X_*(T)^+/\Gamma$，若 $\mu'\in\mu$ 且 $\mathrm{Stab}_\Gamma(\mu')=\Gamma'$，则 $\Gr^{\le\mu}$ 的“相交上同调复形” $IC_\mu$ 对应
  \[
    r_\mu:=\mathrm{Ind}_{\widehat{G}\rtimes\Gamma'}^{{}^LG}r_{\mu'},
  \]
  其中 $r_{\mu'}\in\Rep_{\overline{\QQ}_\ell}(\widehat{G}\rtimes\Gamma')$ 是最高权 $\mu'$ 的不可约表示。
  \item 若 $\mu$ 极小\yizhu{极小（minuscule）余权指对一切正根 $\alpha$ 有 $\langle\alpha^\vee,\mu\rangle\le 1$；此时 $\Gr^{\le\mu}$ 的几何最单纯，$IC_\mu$ 退化为常数层的平移与 Tate 扭。}{A：Bourbaki, \emph{Lie Groups and Lie Algebras}, Ch. VI}，则
  \[
    IC_\mu=\overline{\QQ}_\ell\bigl(\langle\rho,\mu\rangle\bigr)\bigl[\langle 2\rho,\mu\rangle\bigr],
  \]
  其中 $\rho$ 是 $T$ 的正根之和的一半。
\end{enumerate}
\end{conjecture}

借助命题 3.20，相交上同调复形 $IC_\mu$ 应该能在 $\mathit{Hecke}^{\le\mu}$ 上定义一个仍记作 $IC_\mu$ 的“反常层”（为此需要 $IC_\mu$ 的反常性，使得它对 étale 下降数据满足下降）。

最后，这一对
\[
  (\mathit{Hecke}^{\le\mu},IC_\mu)
\]
让我们可以定义 $\Bun_G$ 上的层与 $\Bun_G\times\Spa(E)^{\diamond}$ 上的层之间的\emph{上同调对应}。
